% ============================================================================
% Title:        Introduction
% Author:       Nathaniel Thomas
% Affiliation:  Texas A&M University, Department of Nuclear Engineering
% Email:        nathaniel@tamu.edu
% Date Created: 3/7/2026
% Version:      1.0
%
% Description:
%   A .tex source file describing the rest of the document and providing an
%   overview of the purifier.
%
% =============================================================================

\section{Introduction}

The Texas A\&M FLiNaK hydrodlourination purifier is a chemical
purification system
designed to process FLiNaK salt recieved from commercial vendors, such
as Materion \cite{materion} to reduce the concentration of both metallic
and non-metallic impuritied. \par

\subsection{System description}

The FLiNaK purifier is a hydrofluorination system that removes impurities by
introducing hydrogen gas in combination with hydrogen fluoride gas to react with
impurities present within the molten salt to eliminate impurities. \par

To monitor the progress of purificaiton, a residual gas analyzer is attached
to the system. The gas coming off the top of the reactor during purification is
analyzed for \ch{HF}, \ch{H2O}, and \ch{H2} content, enabling the progress of
purification to be monitored. \par

Following purification, the salt is transferred out of the reaction chamber
through a dip tube, which empties into a secondary stainless steel chamber
where the salt is allowed to cool. From there, salt is allowed to solidify.
Once solidified, the container is transferred into an inert argon glovebox
for further processing. \par

The system is not built inside of an inert atmosphere glovebox, but all
components of the system are sealed, which minimizes the potential for
contamination. In addition, the system is built in a fumehood with a
face velocity of \SI{100}{ft/s} at a sash height of \SI{18}{in},
ensuring operators are protected from any exposure to hazardous \ch{HF}.
There is also an HF detector in the fumehood which detects the precense of
\ch{HF}, which will issue an alarm in the presense of HF, minimizing the
chance of exposure. \par

\subsection{Document structure}

This document is a comprehensive overview of the FLiNaK purifier. It includes
details of construction, such as materials used, part numbers of components
and equipment used, as well as the rational behind various design decisions
(\Cref{chap: construction}). It also contains details on materials and cost
(\Cref{chap: cost}), programming used to control the system
(\Cref{chap: programming}), operation instructions (\Cref{chap: operation}),
operational history (\Cref{chap: operational history}),
and important details on the equipment
used for operation (\Cref{chap: equipment}).

\subsection{Literature review}

% TODO: This section goes over the history of flouride salt purification.
% go based off of Kyle's lit review.

\pagebreak
